첫째 줄에 사진의 수 $N$과 부원의 수 $M$이 공백을 사이에 두고 주어진다. ($1 \le N,M \le 100$)

둘째 줄부터 $N$개의 줄에 걸쳐 각 사진에서 부원들이 눈을 뜨고 있는지를 나타내는 길이 $M$의 문자열이 주어진다.

각 문자열의 $j$번째 문자가 \texttt{1}이면 $j$번 부원이 해당 사진에서 눈을 뜨고 있다는 뜻이고, \texttt{0}이면 눈을 감고 있다는 뜻이다. 문자열은 \texttt{0}과 \texttt{1}로만 이루어지며, 문자 사이에 공백은 없다.